\documentclass[11pt]{amsart}
\usepackage{geometry}                % See geometry.pdf to learn the layout options. There are lots.
\geometry{letterpaper}                   % ... or a4paper or a5paper or ... 
%\geometry{landscape}                % Activate for for rotated page geometry
%\usepackage[parfill]{parskip}    % Activate to begin paragraphs with an empty line rather than an indent
\usepackage{graphicx}
\usepackage{amssymb}
\usepackage{epstopdf}
\DeclareGraphicsRule{.tif}{png}{.png}{`convert #1 `dirname #1`/`basename #1 .tif`.png}

\title{PS 4.19}
\author{Enrique Zuñiga Zuñiga}
%\date{}                                           % Activate to display a given date or no date

\begin{document}
\maketitle
%\section{}
\noindent Se tiene información sobre las calificaciones de 6 exámenes de un grupo de 30 alumnos. Los datos sobre estos exámenes se proporcionan de la siguiente manera:\newline\newline
$$CAL1,1 CAL1,2 ...............CAL1,6  $$\newline
$$CAL2,1 CAL2, 2 ..............CAL2,6  $$\newline
$$CAL30,1 CAL30,2 .........CAI_30,6$$  \newline
Donde: CALi,j es una variable de tipo real que representa la calificación que obtuvo el alumno i en el examen j $(1 < i < 30, 1 < j < 6).$\newline\newline
Escriba un diagrama de flujo que permita calcular lo siguiente:\newline
a) El promedio de calificaciones de cada uno de los 6 exámenes.  \newline
b) El promedio de cada alumno.   \newline
c) El tipo (número) de examen que tuvo el mayor promedio de calificación. Escriba también dicho promedio.  \newline

%\subsection{}



\end{document}  